\chapter{Research and Development Environment}
\label{ch:environment}

This chapter presents the computational infrastructure, software
frameworks, pretrained models and data-access tooling used to design,
train and evaluate the system. Only components actually used in the
project are described, at the exact dependency versions pinned for the
project's environment.

\section{Computing Infrastructure}

\subsection{HPC Cluster and the SLURM Scheduler}

All training and evaluation ran on a shared high-performance computing
cluster managed by \textbf{SLURM} (Simple Linux Utility for Resource
Management)~\cite{yoo2003slurm}, an open-source batch scheduler that
allocates compute nodes to queued jobs according to resource requests and
fair-share priority.

Every stage of the pipeline --- training, evaluation, held-out validation,
classifier probing, and small-scale smoke testing --- runs as its own
batch job, each requesting one GPU, 8 CPU cores and 64\,GB of system
memory on the cluster's GPU partition, and differing mainly in its
walltime request --- a parameter that, as Section~\ref{sec:walltime} of
Chapter~\ref{ch:tasks} describes, turned out to matter a great deal.

Two properties of this environment drove most of the engineering
decisions in Chapter~\ref{ch:tasks}:

\begin{itemize}[leftmargin=*]
  \item \textbf{Compute nodes have no internet access.} Any model weight
        or dataset must be downloaded in advance on the login node, and
        every job is configured to fail immediately, rather than hang
        silently, if it ever attempts a network call it should not need.
  \item \textbf{Jobs are non-interactive and are killed at walltime.} A
        job that exceeds its requested time limit is terminated
        unconditionally, whatever it was doing.
\end{itemize}

\subsection{GPU Hardware}

The cluster is heterogeneous, and the GPU model assigned to a job depends
on which node the scheduler allocates. Two devices were observed across
this project's runs, with measured throughput recorded in the job logs:

\begin{table}[h]
\centering
\caption[GPU hardware and measured runtimes]{GPU hardware observed across
project runs, with measured wall-clock cost of the two main pipeline
halves.}
\label{tab:hardware}
\begin{tabular}{@{}lrrr@{}}
\toprule
\textbf{Device} & \textbf{VRAM} & \textbf{PPO training} & \textbf{Evaluation} \\
\midrule
Tesla P100-PCIE   & 16\,GB & $\sim$20\,h   & $>$4\,h (unfinished) \\
Quadro RTX 5000   & 16\,GB & $\sim$12.5\,h & $\sim$2.5\,h \\
\bottomrule
\end{tabular}
\end{table}

The generator's own memory footprint is modest --- the VAE and UNet
combined occupy an estimated $1826$\,MB --- so the 16\,GB devices
available were never the binding constraint. Wall-clock time was.

\subsection{Environment Provisioning}

A single provisioning step, run once on the login node, creates the
Python virtual environment, installs the pinned dependency set, and ---
most importantly --- \emph{prefetches} every pretrained weight the
compute nodes will later need, including the several TorchXRayVision
DenseNet checkpoints. Dependencies are pinned to exact versions
(\texttt{torch==2.5.1+cu121}, \texttt{diffusers==0.38.0},
\texttt{peft==0.19.1}, \texttt{torchxrayvision==1.5.2},
\texttt{ultralytics==8.4.83}) so that a run months later reproduces the
same numerical environment.

\section{Deep Learning Frameworks}

\subsection{PyTorch}

PyTorch~2.5.1 with CUDA~12.1 is the numerical foundation of the entire
project. Its define-by-run autograd is not merely a convenience here but
a structural requirement: the counterfactual guidance mechanism
differentiates the frozen classifier's loss \emph{with respect to an
intermediate latent inside the diffusion sampling loop}, a gradient path
that must be constructed dynamically at each denoising step. Models run
in half precision on GPU, with selected parameters upcast to full
precision where numerical stability requires it.

\subsection{Diffusers and Stable Diffusion 1.5}

The generator is built on \textbf{Stable Diffusion~1.5}, a latent
diffusion model~\cite{rombach2022} loaded through the Hugging Face
\texttt{diffusers} library. Rather than training a diffusion model from
scratch on a comparatively small medical dataset, the project reuses
SD1.5's pretrained VAE and UNet and adapts only a small parameter subset.

Images of $256\times256$ pixels are encoded by the frozen VAE into a
$4\times32\times32$ latent tensor, an $8\times$ spatial compression that
reduces the cost of every subsequent sampling and guidance step by two
orders of magnitude relative to pixel-space diffusion. Sampling uses
\textbf{DDIM}~\cite{song2021ddim}, whose deterministic, non-Markovian
reverse process permits a partial forward-then-reverse trajectory --- the
mechanism by which an existing radiograph is edited rather than a new
image synthesised from noise.

\subsection{PEFT and LoRA}

Full fine-tuning of the SD1.5 UNet (866~million parameters) was out of
scope for the available compute. \textbf{LoRA}~\cite{hu2021lora}, applied
through the Hugging Face \texttt{peft} library, instead freezes the
pretrained weights and injects trainable low-rank matrix pairs into
selected projections --- here the attention layers' query, key, value and
output projections, at rank 16 with $\alpha=16$ and dropout $0.05$.

The resulting parameter economy is exact: \textbf{6{,}782{,}976 trainable
parameters out of 866{,}303{,}940}, or \textbf{0.78\%} of the model. This
is what made domain adaptation feasible within the internship's compute
budget.

\section{Pretrained Model Zoo}

The system depends on three frozen pretrained networks, none of which is
ever updated by training. They supply, respectively, the guidance
gradient, the spatial constraint, and part of the preservation signal.

\subsection{TorchXRayVision}

\textbf{TorchXRayVision} (XRV)~\cite{cohen2022xrv} is a library of chest
X-ray datasets and pretrained models with a standardised preprocessing
convention. It supplies two distinct components here:

\begin{itemize}[leftmargin=*]
  \item \textbf{Classifiers} --- DenseNet-121 models trained independently
        on different chest X-ray corpora (RSNA, CheXpert, NIH, MIMIC-CXR,
        PadChest) and exposed under a common interface. The availability
        of \emph{multiple independently trained} classifiers over a
        shared label space is what makes the held-out validation of
        Chapter~\ref{ch:results} possible at all; this is the single
        most consequential tooling choice in the project.
  \item \textbf{Segmentation} --- a \textbf{PSPNet} providing per-anatomy
        masks. The lung channels gate where editing is permitted; the
        heart and lung channels together yield the Cardio-Thoracic Ratio
        used as a clinical proxy metric.
\end{itemize}

\subsection{YOLOv8}

\textbf{YOLOv8}~\cite{jocher2023yolov8} (nano variant, loaded via
\texttt{ultralytics}) is used not as a detector but as a \emph{frozen
feature extractor}. Its backbone embedding of an image serves as a
coarse proxy for structural and foreign-object content: penalising drift
in this embedding discourages the generator from destroying tubes, leads
and other hardware. It also underlies the costophrenic-recess proxy
metric.

The report is explicit that this is an approximation. YOLOv8n is
pretrained on COCO, a natural-image corpus, and is not a
radiograph-specific device detector; Section~\ref{sec:limitations}
discusses the consequences.

\subsection{Grad-CAM}

\textbf{Grad-CAM}~\cite{selvaraju2017gradcam} produces class-discriminative
localisation maps by weighting convolutional feature maps with the
gradients flowing into them. A single shared implementation serves two
distinct roles in this project:

\begin{itemize}[leftmargin=*]
  \item \textbf{Diagnostic} --- rendering Grad-CAM overlays to check
        whether the guidance classifier attends to lung pathology or to
        a global brightness/contrast shortcut.
  \item \textbf{Constraint} --- the same map, thresholded and feathered,
        restricts the region in which the PPO agent is permitted to edit
        (Chapter~\ref{ch:tasks}, Section~\ref{sec:gradcam_mask}).
\end{itemize}

\section{Data Access}

The \textbf{RSNA Pneumonia Detection Challenge} dataset~\cite{rsna2018}
is retrieved programmatically through \texttt{kagglehub}, which handles
authentication and local caching so that a compute node can later read
the data with no network access. Images arrive as DICOM and are decoded
with \texttt{pydicom} (with \texttt{pylibjpeg} for JPEG-compressed
transfer syntaxes), min--max normalised to $[0,1]$ and resized to
$256\times256$.

The corpus actually loaded comprises \textbf{26{,}684 studies}, of which
\textbf{6{,}012 are pneumonia-positive} and \textbf{20{,}672 negative} ---
a $22.5\%$ positive rate whose consequences for training are discussed
in Chapter~\ref{ch:tasks}.

\section{Toolchain Summary}

\begin{table}[h]
\centering
\caption{Research and development environment summary.}
\label{tab:toolchain}
\begin{tabular}{@{}p{0.28\textwidth}p{0.62\textwidth}@{}}
\toprule
\textbf{Category} & \textbf{Tools / Technologies} \\
\midrule
Compute infrastructure & SLURM batch scheduler; Tesla P100-PCIE-16GB and
Quadro RTX 5000 GPU nodes; offline (no-internet) compute nodes \\
Numerical framework & PyTorch 2.5.1 + CUDA 12.1, half-precision inference \\
Generative model & Stable Diffusion 1.5 via \texttt{diffusers} 0.38.0; DDIM sampling \\
Parameter-efficient tuning & LoRA via \texttt{peft} 0.19.1 (rank 16, $\alpha$ 16) \\
Pretrained medical models & TorchXRayVision 1.5.2 (DenseNet-121 classifiers, PSPNet segmenter) \\
Object-detection features & Ultralytics YOLOv8n 8.4.83 \\
Explainability & Grad-CAM (custom shared implementation) \\
Data access & \texttt{kagglehub}, \texttt{pydicom}, \texttt{pylibjpeg} \\
Scientific stack & NumPy, SciPy, scikit-image, scikit-learn, Matplotlib \\
Version control & Git \\
\bottomrule
\end{tabular}
\end{table}

\section{Conclusion}

This chapter has described the computational environment of the
internship: a non-interactive, offline, walltime-bounded HPC cluster, a
PyTorch-based generative stack, and a set of frozen pretrained medical
models. The next chapter details how these components were assembled into
a working research pipeline, and how experimental evidence repeatedly
forced that pipeline to be redesigned.
